%% kapa: exam=1 round=29 year=2018 period=2 form=A subject=accounting subject_ko=회계학 seq=28 number=68
%% source: Q-Net 감정평가사 기출문제 (cst00309) artlSeq=5202808 file=extracted/1008/5202808/2018년 제29회 감정평가사 1차시험 기출문제/2018년 제29회 감정평가사 1차시험 기출문제/2교시 A형.pdf
%% conversion: kordoc 4.24.51 -> tools/gen1.py (auto); review=proofread-sonnet; flags=none
\begin{question}{68}
(주)감평은 20x3년도부터 재고자산 평가방법을 선입선출법에서 가중평균법으로 변경하였다. 이러한 회계정책의 변경은 한국채택국제회계기준에서 제시하는 조건을 충족하며, (주)감평은 이러한 변경에 대한 소급효과를 모두 결정할 수 있다. 다음은 (주)감평의 재고자산 평가방법별 기말재고와 선입선출법에 의한 당기순이익이다.
\par{\small\centering
\begin{tabular}{|l|r|r|r|}
\hline
 & \multicolumn{1}{c|}{20x1년} & \multicolumn{1}{c|}{20x2년} & \multicolumn{1}{c|}{20x3년} \\ \hline
기말 재고자산 : & & & \\
\hspace{2em}선입선출법 & ￦1,100 & ￦1,400 & ￦2,000 \\
\hspace{2em}가중평균법 & 1,250 & 1,600 & 1,700 \\ \hline
당기순이익 & ￦21,000 & ￦21,500 & ￦24,000 \\ \hline
\end{tabular}\par}
회계변경 후 20x3년도 당기순이익은? (단, 20x3년도 장부는 마감 전이다.) \par
\begin{choices}
  \item ￦23,500
  \item ￦23,700
  \item ￦24,000
  \item ￦24,300
  \item ￦24,500
\end{choices}
\end{question}
\answer{1}
